%% Disciplined Vibe -- the full seminar deck (one PDF).
%% Master file: preamble + title + agenda, then \input the four unit fragments.
%% Build:  latexmk disciplined_vibe.tex   (lualatex, two passes; see .latexmkrc)
\documentclass[aspectratio=169]{beamer}
\usetheme{tuftedark}
% Unit 3's dependency-graph diagram offsets arrow endpoints with calc syntax
% ($(node)+(dx,dy)$); the theme loads arrows.meta and positioning but not calc.
\usetikzlibrary{calc}

% Drop slide images in images/; include by bare name, e.g.
%   \includegraphics[width=\linewidth]{dino}
\graphicspath{{images/}}

% Placeholder image: \phimg[width-fraction]{label}. Draws a labeled box in place
% of an image not yet made; swap for \includegraphics once the asset exists.
\newcommand{\phimg}[2][0.5]{%
  \tikz\node[draw=tddim, thick, rounded corners, fill=tddim!12,
    minimum width={#1*\linewidth}, minimum height={0.62*#1*\linewidth},
    inner sep=4pt, align=center, text=tddim,
    font=\ttfamily\scriptsize]{\detokenize{#2}};%
}

\title[Disciplined Vibe Coding]{Disciplined Vibe Coding}
\subtitle{The unreasonable effectiveness of language}
\author[Schutera]{Prof.\ Dr.-Ing.\ Mark Schutera}
\institute{\textcolor{DHBWred}{DHBW}\ Ravensburg}
\date{\today}
\setbookversion{\today}{orange disciplined dino}

% Right-edge vertical watermark: a token-id motif instead of the author name.
\renewcommand{\tdwatermark}{[1,\,3,\,6]}

% ---------------------------------------------------------------- citations
% Scientific sources are cited with biblatex under one global numbering for the
% whole deck: every \cite/\slidesources shares a single sequence [1..N], and a
% single shared References package at the end (after unit 4) lists every cited
% paper once, deduplicated. Reuses the notes' shared bib. sorting=none keeps the
% numbers in citation order, so the sources on one slide compress to a tidy
% range, e.g. [1-4]. Only papers are cited; blog posts, talks and courses are
% not.
\usepackage[style=numeric-comp,sortcites,sorting=none,backend=biber]{biblatex}
\addbibresource{../vibecoding_references.bib}
% The References package spans as many slides as it needs (allowframebreaks on
% the frame), so the bibliography no longer has to be crammed onto one frame.
% Size it at \footnotesize (~8pt) to sit with the deck's small body text instead
% of the old 6.5pt one-frame squeeze.
\renewcommand*{\bibfont}{\footnotesize}
\setlength{\bibitemsep}{0.35em plus 0.15em}% a little air between entries now that space is not scarce
\AtEveryBibitem{\clearfield{url}\clearfield{doi}\clearfield{urldate}\clearfield{note}}% drop the full URLs; keep the short arXiv id
\usepackage{multicol}% two-column References frames

% The shared References list is split across four frames (see the References
% frames near the end of the document); at \footnotesize in two columns all the
% cited papers overflow one slide, so we spread them across four so each slide
% keeps a comfortable bottom margin. These categories select each quarter, in
% citation order (6/6/6/5 of the 23 live keys -- the code-review and
% improve-codebase-architecture slides that would add
% martin2008clean/ousterhout2018philosophy/fowler2018refactoring are \iffalse'd
% out in unit4, so those keys are not cited and are omitted here).
% \DeclareBibliographyCategory must live in the preamble. Retag here if the
% cited set changes.
\DeclareBibliographyCategory{refspagea}
\DeclareBibliographyCategory{refspageb}
\DeclareBibliographyCategory{refspagec}
\DeclareBibliographyCategory{refspaged}
\addtocategory{refspagea}{mikolov2013efficient,pennington2014glove,weizenbaum1966eliza,jurafsky2023slp,radford2019gpt2,devlin2018bert}
\addtocategory{refspageb}{raffel2019t5,touvron2023llama,vaswani2017attention,brown2020gpt3,wang2022selfinstruct,ouyang2022instructgpt}
\addtocategory{refspagec}{chung2022flan,nvidia2025gr00t,shukor2025smolvla,kumar2025harnesses,wei2022cot,yao2022react}
\addtocategory{refspaged}{deepseek2024r1,osmani2024seventy,evans2003ddd,beck2002tdd,feathers2004legacy}

% References frame: recolour only the parts that were the red structure colour
% (label, author, location, note, url) to dim grey. The entry title stays white
% (normal text) so it still reads as the primary line.
\setbeamercolor{bibliography entry author}{fg=tddim}
\setbeamercolor{bibliography entry title}{fg=tdfg}
\setbeamercolor{bibliography entry location}{fg=tddim}
\setbeamercolor{bibliography entry note}{fg=tddim}
\setbeamercolor{bibliography item}{fg=tddim}
\DeclareFieldFormat{labelnumberwidth}{\textcolor{tddim}{#1}}

% The side-edge source marker now takes cite keys and prints their numeric
% labels ([1-4]) in the same vertical spot, in place of hand-typed author-years.
\renewcommand{\slidesources}[1]{%
  \begin{tikzpicture}[remember picture,overlay]
    \node[rotate=90,anchor=west,inner sep=0pt,text=tddim,font=\scriptsize]
      at ([xshift=-3mm,yshift=13mm]current page.south east)
      {\cite{#1}};
  \end{tikzpicture}%
}

% ---------------------------------------------------- command-table helpers
% Shared by the opencode command-reference frames across units 2, 3, and 4:
%   \gcmd{/x}  a slash command, red mono (its identity -- what you type).
%   \gkey{k}   a key binding, plain mono.
%   \gpal      reached from the command palette (no slash, no default key).
%   \gsep      thin gap separating a slash from its key on one line.
\newcommand{\gcmd}[1]{\alert{\texttt{#1}}}
\newcommand{\gkey}[1]{\texttt{#1}}
\newcommand{\gpal}{\textcolor{tddim}{palette}}
\newcommand{\gsep}{\hspace{0.7em}}

% A group slide is one booktabs table: "run it" beside "what it does". For a
% key-only action the description opens with its \textbf{name}, since the key,
% not a slash, is its identity.
\newenvironment{cmdtable}{%
  \vskip0.2em\small
  \renewcommand{\arraystretch}{1.1}%
  \begin{tabular}{@{}>{\raggedright\arraybackslash}p{0.26\textwidth}%
                    @{\hspace{5mm}}%
                    >{\raggedright\arraybackslash}p{0.64\textwidth}@{}}%
    \toprule
    \textbf{Run it} & \textbf{What it does}\\
    \midrule}{%
    \bottomrule
  \end{tabular}}

\AtBeginSection[]{\begin{frame}\sectionpage\end{frame}}

\begin{document}

\begin{frame}[plain]
  \titlepage
\end{frame}

% ---------------------------------------------------------------- why
\begin{frame}{Objectives}
  By the end of this lecture series you will have learned,  \\
  \begin{itemize}
    \item The algorithmic foundations and technical intuition of \textbf{large language models} (LLMs).
    \item Put LLMs to use as an \alertbf{engine for vibe coding}.
    \item How to set up a sophisticated \textbf{structure for disciplined} vibe coding.\\
  \end{itemize} 
  \vspace{1.5em}
  The rhythm will be theoretical primers followed by hands-on vibe coding.
\end{frame}


\begin{frame}{Preparing the Workbench}
  \begin{withmargin}
    \begin{tdmain}
      \textbf{We teach on small local models by design.} The hands on sections run a
      \alertbf{small model} served by \textsc{ollama} through the \texttt{opencode}
      terminal client. One bootstrap script installs the whole stack and pulls qwen3:1.7b.
      \vskip0.7em
      \vskip0.9em
      {\footnotesize\color{tddim}There is still value in reading docs as a human:
      \begin{itemize}
        \item \href{https://opencode.ai/docs}{opencode.ai/docs}
        \item \href{https://docs.ollama.com}{docs.ollama.com}
      \end{itemize}}
    \end{tdmain}
    \begin{tdmargin}%
      
      {\small To \alertbf{Set Up} clone the repo and run its bootstrap script. The README walks
      you through macOS, Linux, and Windows.}\par\smallskip
      {\tiny\href{https://github.com/Quillstacks/disciplinedvibe}{github.com/Quillstacks/disciplinedvibe}}
    \end{tdmargin}
  \end{withmargin}
\end{frame}




% ---------------------------------------------------------- audience poll
\begin{frame}{Before we begin, three questions.}
  
  \vskip1em
  \begin{columns}[T,onlytextwidth]
    \begin{column}{0.31\textwidth}
      \textbf{1}\par\smallskip
      Who has access to an \alert{LLM} and uses it regularly?
    \end{column}
    \begin{column}{0.31\textwidth}
      \textbf{2}\par\smallskip
      Who knows how these \alert{systems work}?
    \end{column}
    \begin{column}{0.31\textwidth}
      \textbf{3}\par\smallskip
      Who knows how to \alert{program} in the classical sense?
    \end{column}
  \end{columns}
\end{frame}


% Units share one global citation numbering (see the citations note above): no
% per-unit refsection, and the References package at the end of the deck prints
% every cited paper once.
%% Unit 1 -- An introduction to AI agents.
%% Ported from the standalone "AI Agents" lecture (KIT sdqbeamer deck, L2):
%% same layout and content, retyped against the tuftedark theme API.
%% No preamble here; \input by the master (disciplined_vibe.tex).
%%
%% Scientific sources are cited on the slide they support via \slidesources
%% (cite keys -> a numeric [n] marker up the right edge); the closing frame
%% prints this unit's cited papers with \printbibliography. Biblatex and the
%% per-unit refsection are set up in the master (disciplined_vibe.tex).

\providecommand{\afull}{$\bullet$}
\providecommand{\aemp}{\textcolor{tddim}{$\circ$}}

\sectionsubtitle{History, Theory and Mechanics}
% ======================================================================
\section[Foundations]{Foundations}
% ======================================================================

% ------------------------------------------------ software 1.0 / 2.0 / 3.0
\begin{frame}{The hottest new programming language is English}
  
  \vskip0.8em
  \begin{columns}[T,onlytextwidth]
    \begin{column}{0.31\textwidth}
      \textbf{Software 1.0}\par\smallskip
      {\small Humans write explicit instructions in code, and the machine
      executes them.}
    \end{column}
    \begin{column}{0.31\textwidth}
      \textbf{Software 2.0}\par\smallskip
      {\small Gradient descent writes the weights of neural networks from
      data.}
    \end{column}
    \begin{column}{0.31\textwidth}
      \textbf{Software 3.0}\par\smallskip
      {\small The program is a natural-language prompt, and the LLM
      compiles it to behaviour.}
    \end{column}
  \end{columns}
  \vskip1em
  \begin{tdblock}
    LLMs are a \alertbf{new computing substrate}, with prompts as programs and context windows as RAM.
  \end{tdblock}
\end{frame}


% ------------------------------------------------------- vibe coding: floor/bar
\begin{frame}{Raise the floor, keep the bar}
  \vskip0.8em
  {\small\textbf{Raise the floor.}\; Everyone can now build
  software they could not build before. Fast.}
  \begin{tdblock}
    State your intent, accept the diffs, run the result, embrace the exponentials and iterate
    without necessarily touching or reading every line. That working style is
    \alertbf{vibe coding}.
  \end{tdblock}
  \vspace{1.5em}
  {\small\textbf{Keep the bar.}\; AI engineering means you adjust your workbench to this abstraction layer.}
  \begin{tdblock}
    Most professional work requires rigor. As you will still own
    the code, the quality, and are accountable for whatever you ship.
    That working style is what we term \alertbf{discipline}.
  \end{tdblock}
\end{frame}




% ------------------------------------------------------------- quest: tokens
% \tok -- a single token chip, used to lay a word out as a sequence of tokens.
\newcommand{\tok}[1]{\tikz[baseline]{\node[draw=black!25,rounded corners,
  fill=white,inner sep=3pt,text=black,font=\ttfamily\small,anchor=base]{#1};}}
\begin{frame}{Making the machine understand words.}
  \slidesources{mikolov2013efficient,pennington2014glove,weizenbaum1966eliza,jurafsky2023slp}
  \begin{withmargin}
    \begin{tdmain}
      {\small\color{tddim}Word-level vocabulary}\\[0.4em]
      \tok{banana}\;$\rightarrow$\;\texttt{[24\,907]}\\[0.3em]
      {\footnotesize One token, one entry, one id. A larger vocabulary yields longer
      subwords, representing complex semantics and patterns directly.}
      \vskip1.2em
      {\small\color{tddim}Subword vocabulary}\\[0.4em]
      \tok{ba}\tok{na}\tok{na}\;$\rightarrow$\;\texttt{[3021, 512, 512]}\\[0.3em]
      {\footnotesize Three tokens from two entries, \texttt{ba} and \texttt{na}; the
      repeated \texttt{na} reuses id \texttt{512}. Fewer merges keep the vocabulary
      small and character-based, which improves generalization.}
    \end{tdmain}
    \begin{tdmargin}
      In the beginning was the \textcolor{DHBWred}{token}: a unit of text, a word
      or subword. It is the atomic unit a model reads and writes; a vocabulary is
      the fixed set of tokens it knows. Tokenizing is how text is sliced into
      entries of that set.
      \begin{tdblock}
        LLMs use vocabularies of 30{,}000+ tokens.
      \end{tdblock}
    \end{tdmargin}
  \end{withmargin}
\end{frame}


% ----------------------------------------------------------------- LLM basics
\begin{frame}{Making the machine understand language}
  \sideimage[0.42]{ghibli_parrot}
  \sidecaption{An LLM: the stochastic parrot that predicts the next token.}
  \slidesources{radford2019gpt2,devlin2018bert,raffel2019t5,touvron2023llama}
  \begin{sidebody}
    {\small
    \setlength{\abovedisplayskip}{4pt}\setlength{\belowdisplayskip}{4pt}%
    LLMs learn the structure and semantics of
    language by \alert{predicting the next token},

    \[\begin{aligned}
      \hat{y}_{t+1} &= \theta(x_{1:t}).\\
    \end{aligned}\]


    The free lunch, an implicit ground truth inside every
    token sequence, enables unsuperised learning of $\theta$ on unlabeled data, 

    \[\begin{aligned}
      \mathcal{D} &= \{(x_t, \tilde{y}_t)\}_{t=1}^{N},\\
      \textcolor{DHBWred}{x_{t+1}} &= \textcolor{DHBWred}{\tilde{y}_t}.
    \end{aligned}\]}
  \end{sidebody}
\end{frame}



% ------------------------------------------------------------------ attention
\begin{frame}{Attention is all you need}
  \slidesources{vaswani2017attention}
  \begin{withmargin}
    \begin{tdmain}
      Attention compares token vectors: the closer a query and a key point in the
      same direction, the more that token attends to the other, which reflects \alertbf{semantic similarity}.
      \vskip1.0em
      \centering
      \begin{tikzpicture}[>=Latex,scale=1.15]
        \coordinate (O) at (0,0);
        \coordinate (K) at (2.13,1.49);
        \coordinate (P) at (2.13,0);
        \draw[very thick,tdfg,->] (O) -- (3.2,0) node[right] {$\mathbf{Q}$};
        \draw[very thick,tdfg,->] (O) -- (K) node[above right] {$\mathbf{K}$};
        \draw[tddim] (0.72,0) arc (0:35:0.72);
        \draw[dashed,tddim] (K) -- (P);
        \draw[line width=1.6pt,DHBWred] (O) -- (P);
        \node[DHBWred,font=\footnotesize,below=2pt] at (1.06,0) {$\mathbf{Q}\mathbf{K}^\top$};
      \end{tikzpicture}
      \par\medskip
      {\small\color{tddim}The red projection is the dot product $\mathbf{Q}\mathbf{K}^\top$.}
    \end{tdmain}
    \begin{tdmargin}
      \textbf{Q}, query: what the token is looking for.
      \par\smallskip
      \textbf{K}, key: what a token offers to be matched against.
      \par\smallskip
      \textbf{V}, value: the information a token passes on once matched.
      \par\medskip
      Notice how these computations can be done in \alert{parallel}, 
      and accelerated by GPUs.
    \end{tdmargin}
  \end{withmargin}
\end{frame}



% ------------------------------------------------------------------ context
\begin{frame}{Context is all you have}
  \begin{withmargin}
    \begin{tdmain}
      Each token turns into a query and a key, and attention scores every one
      against every other, a grid of $\mathbf{Q}\mathbf{K}^\top$ interactions.
      This runs token by token over one flat sequence, the \alertbf{context}.
      \vskip1.1em
      \centering
      \begin{tikzpicture}[>=Latex,
          tokcell/.style={draw=tddim,rounded corners=2pt,minimum size=0.55cm,inner sep=0pt},
          faded/.style={draw=tddim!45}]
        \foreach \x in {0.62,1.24,1.86}{\node[tokcell,faded] at (\x,0){};}
        \foreach \x in {2.48,3.10,3.72,4.34,4.96,5.58,6.20}{\node[tokcell] at (\x,0){};}
        \draw[DHBWred,thick,rounded corners=3pt] (2.18,-0.35) rectangle (6.50,0.35);
        \node[DHBWred,font=\scriptsize] at (4.34,0.66) {context window};
        \node[tokcell,draw=DHBWred,fill=DHBWred,text=white,font=\scriptsize]
          (yn) at (7.35,0){$\hat{y}$};
        \draw[->,DHBWred,thick] (6.55,0) -- (yn.west);
      \end{tikzpicture}
      \par\medskip
      {\small\color{tddim}Context is limited - think Goldfish. Its size is capped by GPU memory.}
    \end{tdmain}
    \begin{tdmargin}
      Every token in the window competes for attention:
      \begin{itemize}
        \item your prompt,
        \item system prompts,
        \item the files you shared,
        \item tool and search results,
        \item the conversation so far,
        \item and the model's own replies.
      \end{itemize}
      \par\smallskip
      Irrelevant tokens still enter the $\mathbf{Q}\mathbf{K}^\top$ scores and
      crowd out the signal; doubling the tokens quadruples the compute.
    \end{tdmargin}
  \end{withmargin}
\end{frame}



% ------------------------------------------------------ prompting
\begin{frame}{Still a parrot: prompting}
  \slidesources{brown2020gpt3}
  \begin{withmargin}
    \begin{tdmain}
      \begin{block}{Prompt injection}
        Steer the model by engineering the input prompt, \alertbf{injecting instructions} or
        context that guide and nudge it toward a specific task without retraining.
      \end{block}
      \vskip0.6em
      This is zero-shot, the task lives entirely in the tokens you place in the
      context, so a good prompt is a good program.
    \end{tdmain}
    \begin{tdmargin}
      \textit{``The product is awesome is a good review.}\\
      \textit{The product breaks easily.''}\\
      $\rightarrow$ \textit{[LLM generates the response]}
    \end{tdmargin}
  \end{withmargin}
\end{frame}

% ------------------------------------------------------ instruction tuning
\begin{frame}{Still a parrot: Instruction tuning}
  \slidesources{wang2022selfinstruct,ouyang2022instructgpt,chung2022flan}
  \begin{withmargin}
    \begin{tdmain}
      \begin{block}{Instruction Behavior Training}
        \alertbf{Fine-tune} the model to follow an instruction. At this point the
        \textit{free lunch} is over. An \alert{instruct dataset} pairs each instruction with a label:
        Self-Instruct, InstructGPT, FLAN.
      \end{block}
      \vskip0.6em
      Prompting bends behaviour at inference; fine-tuning writes it into the weights.
      Prompt-to-tune is a recurrent pattern in LLM development.
    \end{tdmain}
    \begin{tdmargin}
      \textit{Instruction: ``Is this review positive?}\\
      \textit{The product breaks easily.''}\\
      $\rightarrow$ \textit{Response: ``No, it is negative.''}      
    \end{tdmargin}
  \end{withmargin}
\end{frame}







% ------------------------------------------------------------- quest: poem
\begin{frame}{\alertbf{Vibe} a poem on bananas}
  \sideimagecover[0.30]{banana_1}
  \sidecaption{The muse for your four-line banana poem.}
  \begin{sidebody}
    If not done yet, \alertbf{set up} your local LLM: clone the repo and run its bootstrap script.\par\smallskip
    {\small\href{https://github.com/Quillstacks/disciplinedvibe}{github.com/Quillstacks/disciplinedvibe}}
    \vspace{0.5em}
    \begin{itemize}
      \item Get comfortable with Open Code
      \item Write at least one poem, what works well, what does not?
      \item Share your best poem in the chat.
    \end{itemize}
    
    \vspace{1.5em}
    Not very useful outside the classroom, but it hints at the \alertbf{semantic capability} of LLMs and their speed, and
    at the applications they open up.   
  \end{sidebody}
  
\end{frame}






% References are collected in one shared package at the end of the deck
% (see disciplined_vibe.tex), not per unit.

%% Unit 2 -- Driving opencode: the commands and keys worth knowing.
%% Commands are grouped by theme, several to a slide, as a two-column booktabs
%% table: how to run it on the left, what it does on the right. Default keys are
%% opencode's out-of-the-box bindings; the leader is ctrl+x.
%% No preamble here; \input by the master (disciplined_vibe.tex).

% The command-table helpers (\gcmd, \gkey, \gpal, \gsep and the cmdtable
% environment) now live in the master preamble (disciplined_vibe.tex), since the
% command tables are shared across units 2, 3, and 4.

\sectionsubtitle{When Language becomes Action}
% ======================================================================
\section[Vibing]{Driving opencode}
% ======================================================================



\iffalse
% -------------------------------------------------------- historical arc
\begin{frame}{A short history of NLP}
  \begin{columns}[T,onlytextwidth]
    \begin{column}{0.24\textwidth}
      \textbf{Pre-2017}\par\smallskip
      {\small N-gram and statistical models, and word embeddings
      such as Word2Vec and GloVe. Early
      chatbots existed, such as ELIZA.}
    \end{column}
    \begin{column}{0.24\textwidth}
      \textbf{2017: the Transformer}\par\smallskip
      {\small The \alert{Transformer} enables highly
      efficient, scalable natural language processing (NLP).}
    \end{column}
    \begin{column}{0.24\textwidth}
      \textbf{2018 to 2021: LLMs}\par\smallskip
      {\small Billions of parameters. Fast progress on pre-training, few-shot
      learning, and RLHF yields BERT, GPT-2, T5, LLaMA.}
    \end{column}
    \begin{column}{0.24\textwidth}
      \textbf{2022, or 0 A.ChatGPT}\par\smallskip
      {\small Mainstream conversational LLMs: GPT-3, and
      from there into vision, language, action models such as GR00T
      and SmolVLA.}
    \end{column}
  \end{columns}
\end{frame}
\fi

% -------------------------------------------------------- tool injection
\begin{frame}{Designing Interfaces by Injection}
  \vskip0.4em
  \begin{withmargin}
    \begin{tdmain}
      The model only reads and writes tokens. A tool still needs an
      interface, so we make one the same way as before: by injecting its
      description into the prompt, then reading the call back out of the reply.
      \vskip0.8em
      \begin{itemize}
        \item The \alertbf{system prompt} specifies the tool, its arguments, and the
          format for calling it.
        \item The model writes that call as ordinary text; a \alertbf{parser} spots
          it and runs the real tool.
        \item The result comes back as tokens and get injected as observations, and the model continues.
      \end{itemize}
    \end{tdmain}
    \begin{tdmargin}
      \begin{tdcodet}{System prompt}
        tool weather(city)\\
        call <tool>...</tool>
      \end{tdcodet}
      \vskip6pt
      \begin{tdcodet}{LLM Tool use}
        <tool>\\
        weather(Ravensburg)\\
        </tool>
      \end{tdcodet}
    \end{tdmargin}
  \end{withmargin}
\end{frame}


% ----------------------------------------------------- what MCP is
\begin{frame}{MCP and plug in outside tools}
  \begin{withmargin}
    \begin{tdmain}
      The \alertbf{Model Context Protocol} is a standard way to connect the agent
      to outside tools and data: a database, a web service, a file store, a program, each
      exposed through a small server.
      \vskip1.0em
      \begin{tdblock}
        MCP is a
        protocol. The agent host connects, discovers them at runtime, and calls
        them. Write to the standard once, and any MCP client can use it.       
      \end{tdblock}
      \vskip0.5em
      {\small }
    \end{tdmain}
    \begin{tdmargin}
      This is the structured answer to agent-ready APIs. The engineering way with proper auth and 
      hardened.
    \end{tdmargin}
  \end{withmargin}
\end{frame}



% ----------------------------------------------------------------- action
\begin{frame}{The one Acting}
  \slidesources{nvidia2025gr00t,shukor2025smolvla}
  \sideimagecover[0.33]{amznghand}
  \begin{sidebody}
    \begin{tdblock}\small
      {\alert{\textbf{Agentic}}, from Latin \textit{agere}, ``to act,'' the
      root of \textit{agent}, one who acts.}
    \end{tdblock}
    \vskip0.15em
    {\footnotesize
    \begin{itemize}\itemsep1pt
      \item Emits its action in a fixed structured format using tags.
      \item Stops at \texttt{</action> or </tool>} once emitted.
      \item A parser then runs the action.
    \end{itemize}}
    \vskip0.15em
    {\footnotesize An action can be many things. Communication and subagent orchestration.
    Information gathering through search, queries,
    and retrieval. It uses tools, from API calls to code execution. Or it acts on
    the environment, driving digital or embodied interfaces.}
  \end{sidebody}
\end{frame}

% ============== The harness that wraps the agent ======================
% With thought, action, and observation on the table, opencode is the harness
% around that agent; sandboxing is the boundary it runs in.

% ----------------------------------------------------- what a harness is
\begin{frame}{The Harness}
  \slidesources{kumar2025harnesses}
  \sideimagecover[0.33]{harness}
  \begin{sidebody}
    The machine got an interface, now it's time for the human to get one to take
    control.
    \vskip0.6em
    \begin{tdblock}
      The \alertbf{harness} is everything wrapped around the raw model to make it
      reliable: the client, the system prompt, the skills, the tests, the
      permission gates, the loop that runs them.
    \end{tdblock}
    \vskip0.6em
    {\texttt{opencode} is your primary harness interface: type
    \alertbf{\texttt{/}} to open the command menu and pick by name.}
    \vskip0.5em
    {\small It has a set of \alertbf{built-in tools} which can be
    called directly: \texttt{bash}, \texttt{read}, \texttt{write},
    \texttt{edit}, \texttt{glob}, \texttt{grep}, \texttt{webfetch}.}
  \end{sidebody}
\end{frame}


% ====================== Getting oriented ==============================
\begin{frame}{Getting started}
  \begin{cmdtable}
    \gcmd{/init} &
      Scans the project and writes an \texttt{AGENTS.md} of your stack, scripts,
      and conventions, so every later session starts oriented and in \alertbf{context}. Run once per repo.\\[4pt]
    \gcmd{/help} &
      Opens the in-app help: the command list, the current keybinds, and where
      the docs live. Your first stop when a key surprises you.\\[4pt]
    \gcmd{/exit}\gsep\gkey{ctrl+x q} &
      Ends the session and drops you back to the shell. \gkey{ctrl+c} on an empty
      prompt quits too.\\
  \end{cmdtable}
\end{frame}



% ====================== Sessions ======================================
\begin{frame}{Steer by Context}
  \begin{cmdtable}
    \gcmd{/compact}\gsep\gkey{ctrl+x c} &
    Summarises the conversation into a short brief and continues from it,
    reclaiming context-window space without losing the thread. Don't wait too long.\\[4pt]
    \gcmd{/new}\gsep\gkey{ctrl+x n} &
      Starts a fresh session with an empty context. Your project stays; the
      conversation so far is dropped. Reach for it when the topic changes.\\[2pt]
    \gcmd{/sessions}\gsep\gkey{ctrl+x l} &
      Lists your earlier sessions and jumps to the one you pick.\\[2pt]
    \gkey{ctrl+r} &
      \textbf{Rename} the current session to a memorable title, so you find it
      again in the list.\\[2pt]
    \gkey{esc} &
      \textbf{Interrupt} the agent mid-run the moment it heads the wrong way,
      instead of waiting for it to finish.\\
  \end{cmdtable}
\end{frame}


% -------------------------------------------- gate vs cage (two layers)
\begin{frame}{Gating and sandboxing}
  \vskip0.3em
  {\small Gating keeps the human on the loop, the sanbox stops agent escapes. 
  Latest when you auto, you want a sandbox.
  }
  \vskip0.8em
  \begin{columns}[T,onlytextwidth]
    \begin{column}{0.47\textwidth}
      \begin{tdblock}
        \small
        \textbf{Gating}\par\smallskip
        \alertbf{opencode asks} before it edits or runs: \texttt{allow},
        \texttt{ask}, \texttt{deny}, set per tool in \texttt{opencode.json}.
        Ships in the box.
        \vskip0.5em
        Stops surprises. It cannot stop a command you already approved.
      \end{tdblock}
    \end{column}
    \begin{column}{0.47\textwidth}
      \begin{tdblock}
        \small
        \textbf{Sandboxing}\par\smallskip
        A container or VM. A project mounted, network restricted, \texttt{opencode}
        running inside it. You build it yourself.
        \vskip0.5em
        \alertbf{Stops escapes}. If something goes wrong it wrecks the box,
        not your machine.
      \end{tdblock}
    \end{column}
  \end{columns}
  \vskip0.9em
  \centering
\end{frame}
















% -------------------------------------------------------------- reasoning
% Two ways to get an inner monologue: elicit it by prompting (CoT, and its
% agentic extension ReAct) on the left, or train it in with RL (R1, o1) right.
\begin{frame}{Reasoning or the Inner Monologue}
  \slidesources{wei2022cot,yao2022react,deepseek2024r1}
  \vskip0.4em
  \begin{columns}[T,onlytextwidth]
    \begin{column}{0.48\textwidth}\small
      \textbf{Prompted at inference time}\par\smallskip
      \textit{reasoning coaxed out by the prompt, no weights changed.}
      \vskip0.5em
      \alertbf{Chain-of-thought} prompting appends a cue like \textit{``Let's think
      step by step,''} steering decoding toward a plan rather than a snap answer,
      so the model breaks the task into sub-tasks.
      \vskip0.5em
      We can interleave those reasoning steps with
      actions and observations so thinking and tool calls advance together, called ReAct.
    \end{column}
    \begin{column}{0.48\textwidth}\small
      \textbf{Trained into the weights}\par\smallskip
      \textit{reasoning trained not nudged.}
      \vskip0.5em
      Models like DeepSeek-R1 and OpenAI's o1 are trained with \alertbf{human feedback reinforcement
      learning} to reason before they answer.
      \vskip0.5em
      \begin{tdblock}
        R1 emits a thinking section between the tokens \texttt{<think>} and
        \texttt{</think>} before its answer.
      \end{tdblock}
    \end{column}
  \end{columns}
\end{frame}


% ====================== Interface and input ===========================
\begin{frame}{Interface and input}
  \begin{cmdtable}
    \gcmd{/editor}\gsep\gkey{ctrl+x e} &
      Opens your \texttt{\$EDITOR} to compose a long message with full editing,
      then sends it back to the prompt.\\[4pt]
    \gkey{ctrl+v} &
      \textbf{Paste} clipboard content into the prompt.\\[4pt]
    \gkey{ctrl+c} &
      \textbf{Clear} the prompt box in one stroke.\\[4pt]
    \gcmd{/details} &
      Toggles a tool call's full input and output, to inspect what
      ran in detail.\\[4pt]
    \gcmd{/thinking} &
      Toggles the model's reasoning \texttt{<think>} blocks, if the model has this capability.\\
  \end{cmdtable}
\end{frame}

% ====================== The surfaces you are driving ==================
% Four awareness frames: what a skill, MCP, harness, and sandbox actually
% are. We use only skills in this course; the other three are named so the
% students know them, and so Unit 3 can treat skills as verbs it invokes.

% (Skills frame moved to Unit 3, placed before the grill-me chain.)





% ------------------------------------------------------------- quest: vibe it
\begin{frame}{\alertbf{Vibe} an infinitely running banana}
  \sideimagecover[0.33]{banana_1}
  \sidecaption{Screenshot of a banana again.}
  \begin{sidebody}
    Now vibe an endless runner in a single \texttt{index.html}. 
    Then ask the LLM to run it in your browser 
    and play, then iterate.
    \begin{itemize}
      \item How easy is it to get one specific thing done?
      \item Watch your context and pay attention to the impact of a saturated context.
      \item Reflect on how productive this feels versus how productive you are.
    \end{itemize}
    \vspace{0.5em}
    \begin{tdblock}
      Feel the \alert{floor rising} already? Lines of code have never been cheaper.
    \end{tdblock}
  \end{sidebody}
\end{frame}

% References are collected in one shared package at the end of the deck
% (see disciplined_vibe.tex), not per unit.
%% Unit 3 -- Plan before you build.  (Draws from Ch5 Planning and Alignment.)
%% Restyled to the sparse, one-idea-per-slide house style of units 1-2:
%% one thesis-shaped title, prose-first bodies, EXACTLY one red per rendered
%% frame (the two cmdtable frames keep their \gcmd slash-commands as identity).
%% Scientific claims are cited via \slidesources; the closing frame prints the
%% unit's papers. No preamble here; \input by the master (disciplined_vibe.tex).
%%
%% ORDER follows the software-engineering process:
%%   A. Motivation      -- why naive vibe coding breaks (four failure modes).
%%   B. Tooling setup   -- model card; agents/subagents; models+agents cmds;
%%                         plan vs build modes; the skill mechanism (AGENTS.md,
%%                         a skill is a file).
%%   C. Plan-arc map    -- the plan catalogue in four columns (align, model, scope + slice, spike).
%%   D. Plan arc, in process order -- askuserquestion (align),
%%                         domain-modeling (name), to-prd (scope),
%%                         to-issues (slice), prototype (spike a design question),
%%                         codebase-design (deep module).
%%   E. Breakout        -- plan the banana, write no code.
%% The build half (handoff + TDD, code-review, diagnosing-bugs,
%% improve-codebase-architecture, write-a-skill) now lives in Unit 4.

\sectionsubtitle{On Humansplaining and LLM Understanding}
% ======================================================================
\section[Alignment]{Plan}
% ======================================================================


% ======================================================================
% A. MOTIVATION
% ======================================================================

% === Vibe coding is brittle: the four failure modes ===
\begin{frame}{Pareto hits the vibe. Hard.}
  \slidesources{osmani2024seventy}

  % TODO (light Pocock spine): Pocock's catalogue exists to fix four failures:
  % misalignment (you and the agent mean different things), verbosity (400 lines
  % you never asked for), non-functional code, and architectural decay. Block 2
  % fixes each with a discipline. Unit 3 takes the first two; Unit 4 the last two.
  % Source: Pocock skills catalogue (pocock2025skills, four failure modes).
  \vskip0.3em
  {\small Naive vibe coding is brittle, then fails. Software engineering
  discipline (and a better model) fixes each mode.}
  \begin{columns}[T,onlytextwidth]
    \begin{column}{0.23\textwidth}
      \begin{tdblock}[equal height group=vibefail]
        \small
        \alertbf{Misalignment.}\; you and the agent mean different things by the
        same word.
      \end{tdblock}
    \end{column}
    \begin{column}{0.23\textwidth}
      \begin{tdblock}[equal height group=vibefail]
        \small
        \alertbf{Verbosity.}\; you ask for one thing and get 400 lines, while
        other functions silently break.
      \end{tdblock}
    \end{column}
    \begin{column}{0.23\textwidth}
      \begin{tdblock}[equal height group=vibefail]
        \small
        \alertbf{Non-functional.}\; the high-score save looks wired, but nothing
        persists.
      \end{tdblock}
    \end{column}
    \begin{column}{0.23\textwidth}
      \begin{tdblock}[equal height group=vibefail]
        \small
        \alertbf{Decay and Debt.}\; each build-loop makes project and module
        complexity harder to manage.
      \end{tdblock}
    \end{column}
  \end{columns}
  \end{frame}


% ======================================================================
% B. TOOLING SETUP
% ======================================================================

  \begin{frame}{Reading a model card}
    \begin{withmargin}
      \begin{tdmain}
        \centering
        \resizebox{\linewidth}{!}{%
        \begin{tikzpicture}[
            node distance=2mm,
            seg/.style={draw=tdfg, thick, rounded corners=2pt,
              inner xsep=7pt, inner ysep=6pt, text=tdfg, font=\ttfamily\large},
            seghi/.style={seg, draw=DHBWred, text=DHBWred},
            sep/.style={text=tddim, font=\ttfamily\large},
            lbl/.style={text=tddim, font=\scriptsize, align=center}]
          \node[seg] (fam) {qwen2.5};
          \node[sep, right=of fam] (c1) {-};
          \node[seg, right=of c1] (dom) {coder};
          \node[sep, right=of dom] (c2) {:};
          \node[seghi, right=of c2] (sz) {7b};
          \node[sep, right=of sz] (c3) {-};
          \node[seg, right=of c3] (pt) {instruct};
          \node[sep, right=of pt] (c4) {-};
          \node[seg, right=of c4] (q) {q4\_K\_M};
          \node[lbl, below=2.5mm of fam] {family\\+ version};
          \node[lbl, below=2.5mm of dom] {domain};
          \node[lbl, text=DHBWred, below=2.5mm of sz]  {parameters};
          \node[lbl, below=2.5mm of pt]  {post-training};
          \node[lbl, below=2.5mm of q] {quantization};
        \end{tikzpicture}}
        \vskip1.1em
        \begin{tdblock}
          {\small Two things the name never carries. The \textbf{context window}
          (\texttt{num\_ctx}) is the token budget the model attends to at once,
          paid for in \textbf{GPU memory}. And whether the model \textbf{reasons}
          on a scratchpad before it answers.}
        \end{tdblock}
        \vskip0.5em
        {\footnotesize\color{tddim}Feel free to switch to more capable models now.}
      \end{tdmain}
      \begin{tdmargin}
        {\color{tdfg}\textbf{qwen2.5}}\; base model and release.\par\smallskip
        {\color{tdfg}\textbf{coder}}\; fine-tuned for a domain, here code.\par\smallskip
        {\color{DHBWred}\textbf{7b}}\; \alert{parameters}: seven billion weights
        in memory (15GB) while the model runs.\par\smallskip
        {\color{tdfg}\textbf{instruct}}\; follows instructions and chats; a
        \texttt{base} model only continues text.\par\smallskip
        {\color{tdfg}\textbf{q4\_K\_M}}\; quantization: about 4 bits per weight,
        shrinking memory to 4.4\,GB.
      \end{tdmargin}
    \end{withmargin}
  \end{frame}

% === Agents and subagents (concept, pairs with the cmdtable below) ===
\begin{frame}{An agent is a prompt, its tools, and a model}
  \begin{withmargin}
    \begin{tdmain}
      \textbf{Plan} and \textbf{build} are the two built-in \textbf{agents}. Each
      bundles a system prompt, a set of allowed tools, and a model. You define
      your own the same way: one markdown file in \texttt{.opencode/agent/}.
      \vskip0.8em
      A \alertbf{subagent} is an agent a primary one spawns to handle a single
      self-contained task in its \textbf{own fresh context window}, then hands
      back a short summary. The side quest never clutters your main session.
      \vskip0.9em
      \begin{tdblock}
        {\ttfamily\small
        > @007 shake it, do not stirr.}
      \end{tdblock}
    \end{tdmain}
    \begin{tdmargin}
      Switch primary agents with \texttt{/agent} or \texttt{Tab}; call a subagent
      by name with \texttt{@}.
      \par\medskip
      Delegation is a tool call. Enter agent orchestration.
    \end{tdmargin}
  \end{withmargin}
\end{frame}

% === Models and agents (cmdtable) ===
\begin{frame}{Models and agents}
  \begin{cmdtable}
    \gcmd{/models}\gsep\gkey{ctrl+x m} &
      Switches the model the agent runs on, from your configured providers.\\[4pt]
    \gkey{f2} &
      \textbf{Recent model:} toggles back to the model you used just before, for
      a quick side-by-side on the same task.\\[4pt]
    \gcmd{/agent}\gsep\gkey{ctrl+x a} &
      Switches the active agent, build versus plan, each with its own tools and
      system prompt.\\[4pt]
    \gkey{tab} &
      \textbf{Cycle agent:} steps to the next agent without opening the list.\\[4pt]
    \gkey{right} &
      \textbf{Cycle subagents:} moves between the subagents a session has spawned,
      to read what each is doing.\\
  \end{cmdtable}
\end{frame}

% -------------------------------------------- plan vs build (two modes)
\begin{frame}{Plan and Build mode in opencode}
  \vskip0.3em
  {\small Plan mode reads and reasons; Build mode writes and runs. Like gating,
  the split mitigates surprises. Toggle with \texttt{Tab}.}
  \vskip0.8em
  \begin{columns}[T,onlytextwidth]
    \begin{column}{0.47\textwidth}
      \begin{tdblock}
        \small
        \textbf{Plan}\par\smallskip
        \alertbf{Read only}. opencode explores, explains, and proposes; it never
        edits files or runs commands.
        \vskip0.5em
        Use it to scope the work and agree on the approach before any change lands.
      \end{tdblock}
    \end{column}
    \begin{column}{0.47\textwidth}
      \begin{tdblock}
        \small
        \textbf{Build}\par\smallskip
        \alertbf{Full access}. opencode edits files and runs tools to carry out
        a plan. Here small models default into action.
        \vskip0.5em
        Use it once the plan is clear to both parties. You can build across
        several sessions.
      \end{tdblock}
    \end{column}
  \end{columns}
  \vskip0.9em
  \centering
\end{frame}


% === A skill is just a file the agent reads ===
\begin{frame}{Agentspanning-Capabilities or Skills}
  \begin{withmargin}
    \begin{tdmain}
      A \alertbf{skill} is a markdown file of instructions and context. Invoke it
      by name; the agent opens that file and follows it, turn by turn. That is the
      whole mechanism.
      \vskip0.6em
      Those instructions do more than steer the model. A skill can call an API,
      run a script, or hand a slice of work to a subagent, so one named file
      becomes the interface to agentic workflow automation.
    \end{tdmain}
    \begin{tdmargin}
      Invoke with /skills, then pick one.
      \par\medskip
      Compact, token-frugal skills win: the whole file is read into context every
      time the skill runs.
    \end{tdmargin}
  \end{withmargin}
\end{frame}

% === write-a-skill ===
\begin{frame}{\texttt{write-a-skill}}
  \begin{withmargin}
    \begin{tdmain}
      Every skill so far taught the agent one job. \texttt{write-a-skill} teaches
      it to author the next one in the same shape, so the toolkit grows by its own
      rules. Name one job, list \textbf{3 to 5 bold-lead steps}, close on a single
      principle, all short enough to read in a breath. If the job needs an ``and'', split it,
       and wrap it in another skill.
      \alertbf{skillception}.
      \vskip0.6em
      \begin{tdblock}
        {\ttfamily\small
        > write-a-skill for "writing skills" \\
        Or wait, I will just nano this myself.}
      \end{tdblock}
    \end{tdmain}
    \begin{tdmargin}
      Writing a skill, is a skill. Of course it is easy to forget, that we can also still do stuff ourselves - manually.
    \end{tdmargin}
  \end{withmargin}
\end{frame}


% ======================================================================
% C. PLAN-ARC MAP
% ======================================================================

% === the plan toolkit: the read-only catalogue this unit walks, one frame ===
% Only the plan arc lives here. The build skills (tdd, code-review,
% diagnosing-bugs, improve-codebase-architecture) and the greenfield on-ramp
% belong to Unit 4; keep them off this frame so it maps only what Unit 3 covers.
\begin{frame}{The plan toolkit, read only}
  \vskip0.3em
  \begin{columns}[T,onlytextwidth]
    \begin{column}{0.23\textwidth}
      \begin{tdblock}
        \scriptsize
        \alertbf{Align}\par\smallskip
        \texttt{askuserquestion}
        \par\medskip
        Reach a design you both hold before a line is written.
      \end{tdblock}
    \end{column}
    \begin{column}{0.23\textwidth}
      \begin{tdblock}
        \scriptsize
        \alertbf{Model}\par\smallskip
        \texttt{domain-modeling}
        \par\medskip
        Fix one word to one meaning, then name the pieces.
      \end{tdblock}
    \end{column}
    \begin{column}{0.23\textwidth}
      \begin{tdblock}
        \scriptsize
        \alertbf{Scope + slice}\par\smallskip
        \texttt{to-prd}\par
        \texttt{to-issues}
        \par\medskip
        Write the brief, then split it into small issues.
      \end{tdblock}
    \end{column}
    \begin{column}{0.23\textwidth}
      \begin{tdblock}
        \scriptsize
        \alertbf{Spike}\par\smallskip
        \texttt{prototype}
        \par\medskip
        Throw one away to settle a single design question.
      \end{tdblock}
    \end{column}
  \end{columns}
\end{frame}


% ======================================================================
% D. PLAN ARC, IN PROCESS ORDER
%    align -> name -> scope -> slice -> spike -> design
% ======================================================================

% === askuserquestion: Align before you name ===
\begin{frame}{\texttt{askuserquestion}}
  \begin{withmargin}
    \begin{tdmain}
      The arc starts with \alertbf{askuserquestion}. It does not write a plan; it
      reaches one. The skill interviews you one question at a time, proposes its
      own answer to each and says why, and walks every branch of the decision tree, until it has
      no questions left.
      \vskip0.6em
      \begin{tdblock}
        {\ttfamily\small
        > run askuserquestion on feature\_brief.md\\
        Q: do power-ups stack? (I'd say no)\\
        > no\\
        Q: does score reset on death?\,\ldots\\
        no open questions. aligned.}
      \end{tdblock}
    \end{tdmain}
    \begin{tdmargin}
      One question at a time, and it stops only when alignment is reached, not
      before. In general following proposals is adviced, until it is not. 
      \par\medskip
      And just because a model has no questions left, does not mean everything is perfectly clarified.
    \end{tdmargin}
  \end{withmargin}
\end{frame}

% === domain-modeling: One word, one meaning, everywhere ===
\begin{frame}{\texttt{domain-modeling}}
  \slidesources{evans2003ddd}
  \begin{withmargin}
    \begin{tdmain}
      \texttt{askuserquestion} settled the vocabulary. You have no
      \texttt{GLOSSARY.md} yet; \texttt{domain-modeling} rereads the chat and code
      and writes one line per term. Now code and tests speak the same words every
      session. \textcolor{DHBWred}{Naming the thing is the alignment.}
      \vskip0.6em
      \begin{tdblock}
        {\ttfamily\small
        > domain-modeling}
      \end{tdblock}
    \end{tdmain}
    \begin{tdmargin}
      \begin{tdcodet}{GLOSSARY.md (after the run)}
        obstacle: box to clear\\
        hitbox: what collides()\\
        duck: lowers the hitbox\\
        ramp: speeds up w/ score
      \end{tdcodet}
      \par\smallskip
      Remember how the substrate of these models is similarity within an embedded feature space.
      So whenever possible go with their wording to increase performance and save context tokens.
    \end{tdmargin}
  \end{withmargin}
\end{frame}

% === to-prd: Write down what you're not building ===
\begin{frame}{\texttt{to-prd}, the product requirements document}
  \begin{withmargin}
    \begin{tdmain}
      Start with letting the agent write this document based on the survey, then step in and adjust, extend, and clean up.
      \vskip0.8em
      \begin{tdblock}
        {\ttfamily\small
        Acceptance lines\\
        - the high score survives a page reload\\
        - the score ticks up each step\\
        Out of scope\\
        - no multiplayer, no leaderboard\\
        - no power-ups this pass}
      \end{tdblock}
    \end{tdmain}
    \begin{tdmargin}
      to-prd fixes verbosity: the aligned chat becomes a short PRD in
      \texttt{issues/PRD.md}. Two parts carry the weight, the out-of-scope list and
      one \textcolor{DHBWred}{acceptance line} per slice, a checkable sentence like
      ``the high score survives a page reload''.
    \end{tdmargin}
  \end{withmargin}
\end{frame}

% === to-issues: Slices you can grab ===
\begin{frame}{\texttt{to-issues}}
  \begin{withmargin}
    \begin{tdmain}
      \texttt{to-issues} splits \texttt{issues/PRD.md} into independently
      \alertbf{grabbable} slices, one markdown file each. Every issue carries its
      acceptance line as its ``how tested''.
      \vskip0.9em
      \begin{tdblock}
        {\ttfamily\small
        issues/01-banana-on-screen.md\\
        \phantom{ }tested: the banana runs across the screen\\[3pt]
        issues/02-persist-high-score.md\\
        \phantom{ }tested: the high score survives a reload\\
        \phantom{ }blocked-by: 01}
      \end{tdblock}
    \end{tdmain}
    \begin{tdmargin}
      Reject a horizontal first slice and re-slice until the first one shows a
      number. Check it actually writes the slice.mds into \texttt{issues/}.
    \end{tdmargin}
  \end{withmargin}
\end{frame}

% === prototype: Spike it, learn it, delete it ===
\begin{frame}{Spike by \texttt{prototype}}
  \begin{withmargin}
    \begin{tdmain}
      What if you can not answer a design question? \texttt{prototype} answers it with a
      \alertbf{throwaway spike}. The answer is the deliverable; the scratch code is
      deleted. 
      \vskip0.9em
      {\small Will this library do the job, is the API fast enough, does the real
      data fit the schema? A 20-line script answers all three. Record ``library
      works, 120\,ms median, dates are strings'', delete the script.}
    \end{tdmain}
    \begin{tdmargin}
      The delete step is mandatory. A spike you keep and build on rots into decay:
      it was never designed to last. It was a standalone accepted cost.
    \end{tdmargin}
  \end{withmargin}
\end{frame}


\iffalse
% === codebase-design ===
\begin{frame}{\texttt{codebase-design}}
  \begin{withmargin}
    \begin{tdmain}
      The outcome is a small set of function signatures,
      names and arguments only, with the bodies left empty. Keep the logic
      separate from anything it draws or displays. Name the one job the module
      does, borrow the shared names from \texttt{GLOSSARY.md}, then write only the
      functions that expose that job and hide the work behind them.
      \vskip0.6em
      \begin{tdblock}
        {\ttfamily\small
        > design the core interface\\
        create(config)\\
        step(state, input)\\
        done(state)\\
        node --test  ok, no display}
      \end{tdblock}
      \vskip0.5em
      {\small A few pure functions carry the whole job. Each takes plain
      arguments and returns plain values; the display layer lives outside and
      calls in.}
    \end{tdmain}
    \begin{tdmargin}
      The test of depth: \alertbf{if the core cannot run under \texttt{node
      --test} with no display, the interface is too wide}. A wide draft leaks
      rendering into the core, thin functions that each need a screen to run.
      Catch it with the headless test, then shrink the interface until only the
      pure functions remain.
    \end{tdmargin}
  \end{withmargin}
\end{frame}
\fi

% ======================================================================
% E. BREAKOUT -- plan the banana, write no code.
% ======================================================================

% === Vibe an infinitely running banana ===
\begin{frame}{\alertbf{Vibe} an infinitely running banana}
  \sideimagecover[0.33]{banana_1}
  \sidecaption{The banana runner, planned this time.}
  \begin{sidebody}
    Same endless runner, now under this discipline. Run the whole planning chain
    on \texttt{feature\_brief.md}, back to back, and write no code.
    \begin{itemize}
      \item \texttt{askuserquestion} until it stops asking.
      \item \texttt{domain-modeling} into \texttt{GLOSSARY.md}.
      \item \texttt{to-prd}: Set up a product requirements document.
      \item \texttt{to-issues}: A first vertical slice that serves as a silver bullet.
    \end{itemize}
    \vskip0.4em
    The deliverable is a \alertbf{structured executable plan}. If the agent starts writing
    code, stop it.
  \end{sidebody}
\end{frame}


% References are collected in one shared package at the end of the deck
% (see disciplined_vibe.tex), not per unit.

%% Unit 4 -- Verify without reading every line.  (Ch6 Implementation, Tests, Review.)
%% Restyled to the sparse, one-idea-per-slide house style of units 1-3: one
%% thesis-shaped title, prose-first bodies, EXACTLY one red per rendered frame
%% (the cmdtable frame keeps its \gcmd slash-commands as identity; a red
%% transcript OUTCOME line counts as that frame's one red). Scientific claims are
%% cited via \slidesources; the closing frame prints the unit's papers. Realized
%% visual-lean: the risk graph and the deep-vs-shallow contrast carry their own
%% frames. No preamble here; \input by the master (disciplined_vibe.tex).
%%
%% ORDER follows the build half of the software-engineering process:
%%   A. Handoff      -- the day shift is done; queue the build.
%%   B. Build skills -- TDD, code-review, diagnosing-bugs,
%%                      improve-codebase-architecture, write-a-skill.
%%   C. Breakout     -- vibe the banana again, now tested.

\sectionsubtitle{aber mit höchster Disziplin.}
% ======================================================================
\section[Building]{Software Engineering Knows How already}
% ======================================================================


% ======================================================================
% A. HANDOFF -- the day shift is done; queue the build.
% ======================================================================

% === Now let the agent build one ===
% A transition slide: one idea, no margin column. The takeaway is a full-width
% centred display (it \vfill-s itself), so we let it carry the frame rather than
% cramming it into a tdmain beside an empty tdmargin.
\begin{frame}{And so we build}
  \vskip0.8em
  That was the \alertbf{day shift}: plan, align, and slice by hand,
  human-in-the-loop. You now hold a backlog of trustworthy slices molded into issues.\\
  \vskip0.8em
  You just finished a product owner workflow, ready to hand-off to the Devs.
  Remember the silverbullet.
\end{frame}


\iffalse
% ======================================================================
% B. BUILD SKILLS
% ======================================================================

% === the build toolkit: the read-only catalogue this unit walks, one frame ===
% Mirrors the plan toolkit map in Unit 3: same four-column read-only layout,
% now for the build half. Only the build skills live here; the plan arc
% (askuserquestion, domain-modeling, to-prd/to-issues, prototype) belongs to
% Unit 3, so keep it off this frame so it maps only what Unit 4 covers.
\begin{frame}{The build toolkit, read only}
  \vskip0.3em
  \begin{columns}[T,onlytextwidth]
    \begin{column}{0.23\textwidth}
      \begin{tdblock}
        \scriptsize
        \alertbf{Test}\par\smallskip
        \texttt{test-driven-development}
        \par\medskip
        Turn each acceptance line into a failing test, then build to green.
      \end{tdblock}
    \end{column}
    \begin{column}{0.23\textwidth}
      \begin{tdblock}
        \scriptsize
        \alertbf{Review}\par\smallskip
        \texttt{code-review}
        \par\medskip
        Read the diff as a stranger against standards and spec.
      \end{tdblock}
    \end{column}
    \begin{column}{0.23\textwidth}
      \begin{tdblock}
        \scriptsize
        \alertbf{Diagnose}\par\smallskip
        \texttt{diagnosing-bugs}
        \par\medskip
        Trace the symptom into the pure core, fix the cause.
      \end{tdblock}
    \end{column}
    \begin{column}{0.23\textwidth}
      \begin{tdblock}
        \scriptsize
        \alertbf{Deepen}\par\smallskip
        \texttt{improve-codebase-architecture}
        \par\medskip
        Pull one clean boundary behind the core, proposal only.
      \end{tdblock}
    \end{column}
  \end{columns}
\end{frame}
\fi

% === test-driven-development ===
\begin{frame}{\texttt{test-driven-development}}
  \slidesources{beck2002tdd}
  \begin{withmargin}
    \begin{tdmain}
      Each acceptance line in \texttt{issues/} becomes a test that fails before it
      passes. The skill drives the red, green, refactor loop one small change at a
      time, running \texttt{node --test} after each. \alertbf{Write the failing test
      before the code that satisfies it}, so each feature earns
      a real acceptance.
      \vskip0.5em
      \begin{tdblock}
        {\ttfamily\small
        > test-driven-development on 01-banana-on-screen\\
        RED: banana not shown, test fails\\
        GREEN: banana renders and runs, test passes}
      \end{tdblock}
    \end{tdmain}
    \begin{tdmargin}
      Fix the code, (never) the test, and watch the runner execute. 
      This concept can obviously be extended to integration and end-to-end tests as well.
    \end{tdmargin}
  \end{withmargin}
\end{frame}

\iffalse
% === code-review ===
\begin{frame}{\texttt{code-review}}
  \slidesources{martin2008clean}
  \begin{withmargin}
    \begin{tdmain}
      Read the diff as if a stranger wrote it: drop all prior context so the agent
      cannot defend its own code, then make two passes.
      \begin{itemize}\itemsep4pt
        \item \textbf{Standards:} check the diff against \texttt{AGENTS.md} and
          \texttt{GLOSSARY.md}.
        \item \textbf{Spec:} check it against the issue acceptance line or the
          \texttt{PRD}.
        \item Keep the core in \texttt{game.js} free of canvas and DOM.
        \item \alertbf{Report the problems, do not rewrite the code.}
      \end{itemize}
    \end{tdmain}
    \begin{column}{0.30\textwidth}
      \begin{tdblock}
        {\ttfamily\small
        > run code-review on the diff\\
        STANDARDS: game.js imports canvas, violates AGENTS.md\\
        SPEC: acceptance line met}
      \end{tdblock}
    \end{column}
  \end{withmargin}
\end{frame}
\fi


% === diagnosing-bugs ===
\begin{frame}{\texttt{diagnosing-bugs}}
  \slidesources{feathers2004legacy}
  \begin{withmargin}
    \begin{tdmain}
      Bugs will show up. Fixing it is mandatory, but then 
      hardening your codebase, by fortifying the fix with tests is best.
      \begin{itemize}\itemsep4pt
        \item Reproduce the bug with a test.
        \item Fix and document the bug.
        \item Confirm the fix by rerunning your test.
        \item \alertbf{Add test to suite.}
      \end{itemize}
    \end{tdmain}
    \begin{column}{0.30\textwidth}
      \begin{tdblock}
        {\ttfamily\small
        > diagnosing-bugs on ``some error message''}
      \end{tdblock}
    \end{column}
  \end{withmargin}
\end{frame}

\iffalse
% === improve-codebase-architecture ===
\begin{frame}{\texttt{improve-codebase-architecture}}
  \slidesources{ousterhout2018philosophy,fowler2018refactoring}
  \begin{withmargin}
    \begin{tdmain}
      Decay creeps in when concerns bleed together, like physics tangled into the
      render loop in \texttt{index.html}. The skill reads \texttt{GLOSSARY.md} and
      scans for mixed concerns.
      \begin{itemize}\itemsep4pt
        \item Pick \alertbf{one clean boundary} to pull into the pure core.
        \item Name what moves behind \texttt{game.js}, proposal only.
        \item Never a sweeping rewrite, never a direct edit.
        \item Keep \texttt{createGame}, \texttt{nextState}, \texttt{collides} a deep
          module you can test.
      \end{itemize}
    \end{tdmain}
    \begin{column}{0.30\textwidth}
      \begin{tdblock}
        {\ttfamily\small
        > run improve-codebase-architecture\\
        physics tangled into render loop in index.html\\
        proposal: move step() behind nextState in game.js}
      \end{tdblock}
    \end{column}
  \end{withmargin}
\end{frame}
\fi


% ======================================================================
% C. BREAKOUT -- vibe the banana again, now tested.
% ======================================================================

% === Vibe an infinitely running banana ===
\begin{frame}{\alertbf{Vibe} an infinitely running banana}
  \sideimagecover[0.33]{banana_1}
  \sidecaption{The banana runner, now building.}
  \begin{sidebody}
    Same banana, now under the discipline of tests. Take one acceptance line
    from the plan and run the whole build chain on it, back to back.
    \begin{itemize}
      \item \texttt{test-driven-development}. Write a test, and a feature against it.
      \item \texttt{diagnosing-bugs}: Encounter a bug, fix it, harden against it.
    \end{itemize}
    \vskip0.4em
    The deliverable is a \alertbf{tested banana core} you trust without reading
    every line. And what remains is taste.
  \end{sidebody}
\end{frame}


% References are collected in one shared package at the end of the deck
% (see disciplined_vibe.tex), not per unit.


% -------------------------------------------------------- shared references
% One References package for the whole deck: every paper cited by any unit,
% under one global numbering, deduplicated. Placed here -- after the units,
% before the closing summary -- so it reads as the deck's single source list and
% the deck still ends on the closing statement.
%
% All the cited papers at \footnotesize in two columns overflow one slide, so we
% split the list across four frames, one \printbibliography per frame filtered to
% its quarter of the keys in citation order (6/6/6/5), which keeps a comfortable
% bottom margin on each. Numbering stays global because all four draw from the
% same bibliography; the category filter only decides which entries each frame
% prints, not how they are counted. multicols + allowframebreaks was rejected: it
% emits a spurious empty leading slide. The refspage{a,b,c,d} categories are
% declared in the preamble (biblatex requires \DeclareBibliographyCategory there).
\begin{frame}{References}
  \begin{multicols}{2}
    \printbibliography[heading=none,category=refspagea]
  \end{multicols}
\end{frame}
\begin{frame}{References}
  \begin{multicols}{2}
    \printbibliography[heading=none,category=refspageb]
  \end{multicols}
\end{frame}
\begin{frame}{References}
  \begin{multicols}{2}
    \printbibliography[heading=none,category=refspagec]
  \end{multicols}
\end{frame}
\begin{frame}{References}
  \begin{multicols}{2}
    \printbibliography[heading=none,category=refspaged]
  \end{multicols}
\end{frame}

















% ======================================================================

\begin{frame}{What you learned}
  \begin{columns}[T,onlytextwidth]
    \begin{column}{0.30\textwidth}
      \textbf{Foundations}\par\smallskip
      {\color{tddim}You understand the mechanics of LLMs and have a technical intuition about them.}
    \end{column}
    \begin{column}{0.30\textwidth}
      \textbf{Vibe}\par\smallskip
      {\color{tddim}You are able to run an LLM as an engine for vibe coding, and know how to steer it.}
    \end{column}
    \begin{column}{0.30\textwidth}
      \textbf{Discipline} \par\smallskip
      {\color{tddim}You have the means to make your vibing dependable and aligned.}
    \end{column}
  \end{columns}
\end{frame}

\begin{frame}[plain]
  \vfill
  \centering
  {\LARGE\bfseries Give in to the \alertbf{vibes}, keep the \alertbf{discipline}}\par\vskip1.2em
  {\small Prof.\ Dr.-Ing.\ Mark Schutera \;\textbar\; DHBW Ravensburg}
  \vfill
\end{frame}



% Appendix: the ten skills verbatim, so students can copy them into skills/.
% No refsection -- it cites no papers.
%%% Appendix -- the skills, verbatim, so students can copy them into skills/.
%% One frame per skill, grouped Plan (Unit 3) then Build (Unit 4). Each frame
%% reproduces skills/<name>.md line for line in a tdcodet box (the same dark
%% mono box the GLOSSARY.md frame uses), so what is on the slide IS the file.
%% Source of truth: setup/sandbox/skills/*.md. If a skill file changes, change
%% the matching frame here. No preamble; \input by the master after unit4.

% ======================================================================
\section[Skills]{The Skills, to Copy}
% ======================================================================

% -------------------------------------------------- how to use this appendix
\begin{frame}{Copy these into \texttt{skills/}}
  \begin{withmargin}
    \begin{tdmain}
      Every skill is just a markdown file the agent reads when you name it.
      Nothing here is magic. To get the whole toolkit, create one file per skill
      under \texttt{skills/} and paste the text on the next slides into it,
      keeping the \texttt{name:} and \texttt{description:} header as shown.
      \vskip0.8em
      \begin{tdblock}
        {\ttfamily\small
        skills/grill-me.md\\
        skills/domain-modeling.md\\
        skills/to-prd.md \;\ldots{} one file each}
      \end{tdblock}
    \end{tdmain}
    \begin{tdmargin}
      \alertbf{Plan} skills come first (Unit 3), then \alertbf{Build} skills
      (Unit 4).
      \par\medskip
      The filename must match the \texttt{name:} line; that is the word you
      invoke it by.
    \end{tdmargin}
  \end{withmargin}
\end{frame}

% ============================== PLAN SKILLS (Unit 3) ===================

% -------------------------------------------------------------- grill-me
% Longest skill: shrink to \tiny and let the box wrap each rule naturally so
% every line survives above the footer.
\begin{frame}{Plan skill: \texttt{grill-me}}
  \begin{tdcodet}{skills/grill-me.md}
    \tiny
    name: grill-me\\
    description: Interview the user relentlessly about a plan until you reach a shared understanding.\\[3pt]
    Interview me relentlessly about every aspect of this plan until we reach a shared understanding. The goal is not a document, it is alignment: a design concept we both hold.\\[3pt]
    Rules:\\[2pt]
    1. First, explore the codebase enough to ask informed questions. If a question can be answered by reading the code, read it instead of asking.\\
    2. Ask questions one at a time. Wait for my answer before the next one.\\
    3. For every question, propose your own recommended answer and say why.\\
    4. Walk down each branch of the decision tree, resolving dependencies one by one. Surface the awkward questions a non-expert would forget (edge cases, retroactive data, failure behaviour, out-of-scope boundaries).\\
    5. Do not write a plan or any code yet. Keep going until you have no open questions, then say so.\\[3pt]
    Stop when alignment is reached, not before.
  \end{tdcodet}
\end{frame}

% -------------------------------------------------------- domain-modeling
\begin{frame}{Plan skill: \texttt{domain-modeling}}
  \begin{tdcodet}{skills/domain-modeling.md}
    name: domain-modeling\\
    description: Build the shared vocabulary for the game and record it in GLOSSARY.md.\\[3pt]
    Pin down the words we use for the game so the code and tests all say the same thing.\\[3pt]
    1. List. Collect the nouns and verbs from the chat and the code: obstacle, hitbox, duck, ramp, spawn rate, run length.\\
    2. Define. Give each term one plain line: what it means and where it lives (createGame, nextState, collides).\\
    3. Record. Write these lines into GLOSSARY.md, one term per line.\\
    4. Fix. Find synonyms (two words, one thing) and clashes (one word, two things), then pick one winner and drop the rest.\\[3pt]
    One word, one meaning, everywhere.
  \end{tdcodet}
\end{frame}

% ---------------------------------------------------------------- to-prd
\begin{frame}{Plan skill: \texttt{to-prd}}
  \begin{tdcodet}{skills/to-prd.md}
    name: to-prd\\
    description: Turn an aligned conversation into a product requirements document.\\[3pt]
    Summarise the design concept we reached into a Product Requirements Document. Write it to issues/PRD.md.\\[3pt]
    Structure:\\[2pt]
    - Problem --- the problem the user faces, in one short paragraph.\\
    - Solution --- the approach, in one short paragraph.\\
    - User stories --- a short list of "as a <role>, I want <goal>, so that <reason>".\\
    - Implementation decisions --- the concrete choices we agreed on.\\
    - Testing decisions --- how each story will be verified.\\
    - Out of scope --- what we explicitly decided NOT to do, and why. This is the definition of done.\\[3pt]
    Do not write implementation code. The PRD records the destination, not the route.
  \end{tdcodet}
\end{frame}

% -------------------------------------------------------------- to-issues
\begin{frame}{Plan skill: \texttt{to-issues}}
  \begin{tdcodet}{skills/to-issues.md}
    \tiny
    name: to-issues\\
    description: Slice a PRD into independently grabbable vertical-slice issues.\\[3pt]
    Break the PRD in issues/PRD.md into independently grabbable issues, written as local markdown files in issues/ (one file per issue, e.g. issues/01-*.md).\\[3pt]
    Rules:\\[2pt]
    1. Slice vertically, not horizontally. A good slice cuts through every layer (data, service, dashboard) and produces something visible and testable. A bad slice is "all the schema" or "all the dashboard".\\
    2. The first slice must be a tracer bullet: the thinnest end-to-end path that proves the whole flow works.\\
    3. Record blocking relationships between issues (which must be done before which), so independent issues could be worked in parallel.\\
    4. Each issue states: title, the user-facing outcome, and how it will be tested.\\[3pt]
    If your first proposed slice is horizontal, reject it and re-slice before writing the files.
  \end{tdcodet}
\end{frame}

% -------------------------------------------------------------- prototype
\begin{frame}{Plan skill: \texttt{prototype}}
  \begin{tdcodet}{skills/prototype.md}
    name: prototype\\
    description: Build a throwaway prototype to answer one design question, then delete it.\\[3pt]
    Answer one design question with a throwaway prototype, then delete it.\\[3pt]
    1. Ask. Write the one design question in a single sentence (e.g. tick-based core or event-based?). If you cannot fit it on one line, stop and narrow it.\\
    2. Scratch. Build the smallest scratch script that answers only that question. Skip edge cases, node --test, and GLOSSARY.md vocabulary.\\
    3. Record. Write the answer as one line in the matching issues/ file, or in GLOSSARY.md if it changes the shared model.\\
    4. Delete. Remove the scratch files so no throwaway code survives.\\[3pt]
    The answer is the deliverable, not the code.
  \end{tdcodet}
\end{frame}

% ============================== BUILD SKILLS (Unit 4) ==================

% ----------------------------------------------------------------- tdd
\begin{frame}{Build skill: \texttt{tdd}}
  \begin{tdcodet}{skills/tdd.md}
    name: tdd\\
    description: Implement a change with the red, green, refactor loop.\\[3pt]
    Implement the task using test-driven development.\\[3pt]
    1. Red. Write or locate a single failing test that specifies the next small\\
    \phantom{1. }behaviour. Run node --test and confirm it fails for the right reason.\\
    2. Green. Write the minimum code to make that test pass. Run node --test.\\
    3. Refactor. With the bar green, tidy the code without changing behaviour.\\
    \phantom{3. }Run node --test again.\\[3pt]
    Make one small change at a time and run the tests after each. Do not edit the\\
    test to fit broken code; fix the code. When done, show the diff for review.
  \end{tdcodet}
\end{frame}

% -------------------------------------------------------- codebase-design
\begin{frame}{Build skill: \texttt{codebase-design}}
  \begin{tdcodet}{skills/codebase-design.md}
    name: codebase-design\\
    description: Design a deep module: a lot of behaviour behind a small, simple interface.\\[3pt]
    Design a deep module: hide a lot of behaviour behind a small, simple interface.\\[3pt]
    1. Name one job. Write one sentence for what this module does. One job only.\\
    2. Check GLOSSARY.md. Reuse the shared names from GLOSSARY.md so the interface fits the domain.\\
    3. Draft the small interface. Write only the signatures that hide the work, like createGame, nextState, collides. Nothing extra.\\
    4. Prove it is testable. Confirm each signature runs under node --test with no canvas and no DOM. If not, shrink it.\\
    5. Implement behind it. Now write the code inside those functions. Keep the interface fixed.\\[3pt]
    A deep module hides a lot behind a little.
  \end{tdcodet}
\end{frame}

% ------------------------------------------------------------- code-review
\begin{frame}{Build skill: \texttt{code-review}}
  \begin{tdcodet}{skills/code-review.md}
    name: code-review\\
    description: Review a diff twice, once for standards and once for spec, in a fresh context.\\[3pt]
    Read the diff fresh and report what is wrong.\\[3pt]
    1. Clear. Drop all prior context first, so you read the diff with no memory of writing it.\\
    2. Standards. Check the diff against the rules in GLOSSARY.md and AGENTS.md. The core in game.js must stay free of canvas and DOM.\\
    3. Spec. Check the diff does what the issue acceptance line or the PRD asked, no more and no less.\\
    4. Report. List every problem you found in one message.\\[3pt]
    Report the problems, do not rewrite the code.
  \end{tdcodet}
\end{frame}

% --------------------------------------------------------- diagnosing-bugs
\begin{frame}{Build skill: \texttt{diagnosing-bugs}}
  \begin{tdcodet}{skills/diagnosing-bugs.md}
    name: diagnosing-bugs\\
    description: Find and fix the root cause of a bug, not the symptom.\\[3pt]
    Trace a bug to its root cause in the pure core (createGame, nextState, collides) and fix that, not the surface symptom.\\[3pt]
    1. Reproduce. Write a failing test that triggers the bug. Run node --test and watch it fail.\\
    2. Minimise. Cut the test to the smallest input that still fails. Name what you see using GLOSSARY.md.\\
    3. Find the cause. State the likely cause in one sentence. Add a log to confirm it before you change anything.\\
    4. Fix. Change the cause, remove the log, and run node --test until green.\\[3pt]
    A patch that hides the cause creates two new bugs.
  \end{tdcodet}
\end{frame}

% ------------------------------------------- improve-codebase-architecture
\begin{frame}{Build skill: \texttt{improve-codebase-architecture}}
  \begin{tdcodet}{skills/improve-codebase-architecture.md}
    name: improve-codebase-architecture\\
    description: Find a tangled cluster and propose deepening it into one testable module.\\[3pt]
    Find one tangled cluster and propose deepening it into a single testable module.\\[3pt]
    1. Read. Open GLOSSARY.md for the shared vocabulary before you look at code.\\
    2. Scan. Find mixed concerns, like physics tangled into the render loop in index.html.\\
    3. Cut. Pick one clean boundary to pull into the pure core (createGame, nextState, collides).\\
    4. Propose. Name what moves behind game.js and its interface. Write the proposal only, do not edit code.\\[3pt]
    A boundary you can wrap a node --test around is the goal.
  \end{tdcodet}
\end{frame}

\end{document}
